% >>> À RÉÉCRIRE TOI-MÊME (priorité HAUTE pour l'authenticité) :
% Personnalise : cite une chose CONCRÈTE que ton encadrant t'a apportée (une
% remarque qui a changé ta trajectoire, p. ex. le recentrage sur la mesure
% plutôt que la performance), un membre du jury, un proche. Les détails vécus
% sont ce qui distingue le plus un texte humain d'un texte généré.
\chapter*{Remerciements}
\addcontentsline{toc}{chapter}{Remerciements}
\thispagestyle{plain}

J'adresse mes remerciements à mon encadrant, le \textbf{Professeur DABOUNOU}, pour son encadrement précieux.

Je remercie le corps enseignant de la filière \textbf{Génie Informatique} pour la formation reçue, ainsi que les membres du jury qui ont accepté d'évaluer ce travail.

Ma reconnaissance va enfin à ma famille et à mes camarades, dont les encouragements ne m'ont jamais manqué.
